% This is a template for your written document.
%
% To compile using latexmk on the command line, run the following: 
% latexmk -pdf main.tex

\documentclass[12pt]{article}
\usepackage{setspace}
\singlespace
\usepackage[left=1in,right=1in,top=1in,bottom=1in]{geometry}
\usepackage{graphicx}

\title{\textbf{Bridging Understanding to Speaking: A Language Learning Application for Heritage Speakers}}
\author{Mercy Yankey}
\date{April 8, 2026}

\begin{document}

\maketitle

\section{Introduction}

   Language learning applications have become one of the most widely used tools for acquiring new languages. However, most existing applications are designed with second language learners in mind, often assuming little to no prior exposure to the users target language. The design approach of only working with reasearch on second learners learning style overlooks a  population of learners known as heritage speakers. Heritage Speakers are ndividuals who have been exposed to a language through family or community but may not have developed full proficiency.

	

   Heritage speakers occupy a unique position in language learning. Unlike traditional second language learners, they often possess strong listening and comprehension skills, but may struggle with reading, writing, or formal grammar. As a result, conventional language learning tools, which typically emphasize vocabulary memorization and basic grammar structures, may not effectively address their needs.


	 
   This project explores whether a mobile language learning application can be designed specifically to support heritage speakers, with a focus on a Native Ghanian language known as Twi. The application incorporates features such as flashcards, spaced repetition, production-based prompts, and interactive quizzes to better align with the strengths and challenges of heritage learners.


   The central research question guiding this work is: \textit{Can a language learning application designed with heritage speakers in mind better support their language development compared to traditional approaches?} To investigate this, the project combines insights from heritage language acquisition research with practical application development.


	 This paper presents the motivation behind the project, reviews relevant research on heritage language learning, and combines the design of the application within existing literature. By doing so, it aims to highlight gaps in current language learning tools and propose a more targeted approach for heritage speakers.

\section{Background}

  Heritage language speakers are typically defined as individuals who grow up in households where a language other than the dominant societal language is spoken, but do not fully develop proficiency in it. This can look very different from person to person. Some heritage speakers can understand almost everything but struggle to respond, while others can speak comfortably but lack reading or writing skills.

    A key idea in heritage language learning is that exposure does not equal mastery. Most heritage speakers learn the language informally, through conversation at home, listening to relatives, or interacting in community spaces. Because of this, their language knowledge is often imbalanced. For example, someone might be able to follow a full conversation but not know how to structure a grammatically correct sentence to respond or in writing.
   
   This imbalance is often explained through what researchers call incomplete acquisition. The idea is that the language was learned early, but not fully developed, usually because the dominant language, like English in the U.S., takes over in school and daily life. Over time, this leads to gaps, especially in formal grammar, vocabulary, and production.

	Another important point is that heritage speakers are not the same as second language learners. Second language learners usually start from zero and build up skills in a structured way. Heritage speakers, on the other hand, already have a foundation, but it is inconsistent. With heritage learners style of inconsistencey in learning, traditional learning tools do not always work well for them.

	This matters for this project because most language learning apps are designed with second language learners in mind. They focus heavily on vocabulary memorization and basic grammar drills, which may not match what heritage speakers actually need. Instead, heritage learners may benefit more from tools that focus on production, reinforcement, and filling in specific gaps in their knowledge. Understanding this difference is important for designing more effective learning tools. In this project, it directly motivates the design of an application that prioritizes active use of the language, rather than just recognition or memorization.

  \section{Related Work}

\subsection{Heritage Language Learning}

Research on heritage language learning has shown that heritage speakers do not learn in the same way as second language learners. Instead of starting from zero, they already have some level of exposure to the language, but that knowledge is often incomplete or uneven. This creates a different learning profile that is not always supported by traditional language learning methods.

One of the most common findings in this area is that heritage speakers tend to have stronger listening and comprehension skills, but weaker production abilities. In other words, they are often able to understand the language much better than they can actively use it. This gap becomes more noticeable in structured situations, such as writing or speaking formally, where they may struggle with grammar or vocabulary. Research on heritage speakers’ vocabulary development supports this pattern, showing that productive language skills are often weaker than receptive skills \cite{gharibi2017factors}.

Another key concept in this research is incomplete acquisition. This refers to situations where a speaker begins learning a language early in life, but does not fully develop all aspects of it. This is often due to a shift toward the dominant language, especially in school environments. Over time, this results in gaps in grammar, vocabulary, and confidence when using the language. Studies on heritage language development suggest that early exposure alone does not guarantee full proficiency, and that language abilities may stabilize before reaching native-like competence \cite{montrul2008gender}.

At the same time, it is important to recognize that heritage speakers are not all the same. Their abilities can vary depending on how often they use the language, who they use it with, and whether they have had any formal instruction. Differences in language exposure and use play a major role in shaping these outcomes. For example, research on bilingual development shows that both the quantity and quality of language input significantly influence language proficiency over time \cite{hoff2018bilingual}. Because of this, a single learning approach is often not effective, and more flexible or targeted approaches are needed.

This section connects directly to my project because it highlights why existing language learning tools may not be effective for heritage speakers. If most tools are designed for learners starting from scratch, they may not address the specific gaps that heritage speakers experience. In particular, the consistent gap between comprehension and production suggests that tools should focus more on active language use rather than passive recognition.

\subsection{Technology in Language Learning}

Technology has become a central component of modern language learning, particularly through mobile applications and computer-assisted language learning systems. Early research on Mobile Assisted Language Learning is focused on how digital tools could supplement traditional instruction by providing additional practice, flexibility, and access to learning materials. Levy describes how technology has evolved from simple instructional support to an integrated part of language learning environments, shaping how learners interact with content and practice language skills \cite{levy2009technology}.

More recent work has shifted toward evaluating the effectiveness of mobile language learning applications. Studies comparing platforms such as Duolingo, Babbel, and Busuu suggest that while these tools are widely used and accessible, their effectiveness depends heavily on how well their design aligns with actual learning needs \cite{muckenhumer2023efficacy}. Many applications emphasize repetition, vocabulary drills, and gamified interactions, which can support memorization but may not always translate into strong production skills.

In addition to structured lessons, some research has explored the use of interactive and game-based learning environments. Heritage language development is also shaped by broader educational and social factors \cite{byrnes2005heritage}. For example, projects incorporating storytelling, games, and interactive systems have shown that engagement and contextual learning can improve motivation and retention \cite{papadopoulos2013aladdin}. These approaches highlight the importance of designing learning experiences that go beyond isolated vocabulary practice and instead encourage active language use.

Despite these advancements, most language learning technologies are still designed with second language learners in mind. They typically assume that users are starting from scratch and follow a linear progression of lessons. This design does not account for heritage speakers, who already possess partial knowledge of the language but require support in developing production skills and filling specific gaps. As a result, there is a disconnect between existing research on heritage language learning and the design of current language learning applications.

\subsection{Gaps in Existing Research}

Overall, existing research highlights two key gaps. First, heritage language learners have distinct learning needs that are not fully addressed by traditional second language instruction models. While prior research clearly identifies differences in comprehension, production, and language development, these insights are not always reflected in the design of language learning tools.

Second, while language learning technologies continue to expand, they are primarily designed for learners without prior exposure to the language. This creates a mismatch between how applications are structured and what heritage speakers actually need. In particular, many applications focus on recognition-based tasks rather than encouraging consistent language production, which is a critical area for heritage learners.

Additionally, many languages spoken by heritage communities, particularly African and other underrepresented languages, receive little attention in mainstream language learning platforms. This lack of representation limits access to learning resources and can make it more difficult for heritage speakers to maintain or improve their language skills.

Together, these gaps suggest a need for more targeted systems that account for both the linguistic background of heritage speakers and the limitations of current application design. This project builds on these findings by exploring how a language learning application can be designed specifically to support heritage speakers, with a focus on improving active language use.
\section{Design Rationale}

The design of this application is based on aligning system features with the known challenges faced by heritage language learners. Rather than simply describing what the system does, this section explains why these design choices were made and how they address gaps identified in prior research.

One of the primary design decisions was to prioritize language production over recognition. As discussed in prior research, heritage speakers often demonstrate strong comprehension skills but struggle with actively producing the language. Many existing applications emphasize vocabulary recognition and repetition, which may reinforce passive understanding but do not necessarily improve speaking ability. In response to this, the application includes production-based prompts and response activities that require users to generate language. This approach is intended to strengthen active recall and support more natural language use over time.

Another key design choice is the use of spaced repetition and structured review. While spaced repetition is commonly used in language learning applications, it is particularly relevant for heritage speakers who may have gaps in vocabulary or inconsistent knowledge. By revisiting words over time, the system reinforces weaker areas while still building on existing knowledge. This allows learners to strengthen areas where their understanding may be incomplete without restarting from a beginner level.

The inclusion of multiple interaction modes, such as quick review sessions and game-based activities, is also intentional. Mobile applications are often used in short, informal contexts, so the system is designed to support both focused learning and low-effort engagement. This flexibility reflects how users actually interact with language learning tools and encourages more consistent use of the application.

In addition, this project focuses on the Twi language, which is currently underrepresented in mainstream language learning platforms. Many widely used applications prioritize globally dominant languages, leaving fewer resources available for African and other heritage languages. By designing specifically for Twi, this project addresses both a learning gap and a representation gap, highlighting how technology can support languages that are often overlooked in digital spaces.

Overall, the design of the application reflects a shift away from traditional second language learning models and toward a more targeted approach. By focusing on production, reinforcement, and flexibility, the system aims to better support the unique learning profiles of heritage speakers while also contributing to broader efforts in language preservation and accessibility.

% --- NOTES FOR CONTINUATION FROM WRITING CENTER ---
% could connect more directly to specific features (games, social, translate)
% could reference specific papers again (tie design decisions back to research)
% could expand on underrepresented languages if needed

\section{Conclusion}
This project examined the gap between existing language learning technologies and the specific needs of heritage speakers. While current applications are widely used and accessible, they are primarily designed for second language learners, a  and do not account for the uneven skill development that is common among heritage users.

By focusing on production-based learning, flexible interaction modes, and targeted reinforcement, this work proposes an alternative approach to language learning application design. The development of a Twi language learning application serves as a case study for how technology can be adapted to better support both heritage speakers and underrepresented languages.

Although this draft focuses on the design of the system, future work will involve evaluating its effectiveness through user testing and measuring improvements in language production and engagement. More broadly, this project contributes to ongoing discussions about accessibility, representation, and the role of technology in supporting language learning and preservation.


\bibliographystyle{acm}
\bibliography{bibliography}

\end{document}
